\documentclass[UTF8,11pt]{ctexart}
\usepackage[letterpaper,margin=1in]{geometry}
\usepackage{xcolor,enumitem,booktabs,longtable,tabularx,array,hyperref,amssymb}
\definecolor{aidinggreen}{HTML}{087A4B}
\definecolor{aidingblue}{HTML}{2E74B5}
\hypersetup{colorlinks=true,linkcolor=aidinggreen,urlcolor=aidingblue}
\setlist[itemize]{leftmargin=1.7em,itemsep=0.25em}
\setlist[enumerate]{leftmargin=1.8em,itemsep=0.25em}
\newcommand{\placeholder}[1]{\textcolor{aidinggreen}{[#1]}}
\newcommand{\checkbox}{\ensuremath{\square}}
\newcommand{\sectionrule}{\par\noindent\textcolor{aidinggreen}{\rule{\linewidth}{0.7pt}}\par}

\title{\textcolor{aidinggreen}{AIDING 中英文论文全流程工作规范}}
\author{AIDING Team}
\date{通用团队 SOP}
\begin{document}
\maketitle
\sectionrule
\textbf{最高优先级：}目标期刊最新作者指南、投稿系统字段和编辑部通知始终高于本通用 SOP。所有最终提交均需团队复核并向刘老师确认。
\section{流程总览}
\begin{itemize}
\item 撰写与定稿：中文母稿、双语分支、期刊预演和线下校稿。
\item 首次投稿：Cover Letter、版本锁定、生成稿核对和归档。
\item 返修与回复：意见编号、双语工作稿、逐点回复和位置映射。
\item 录用与校样：Proof、出版问题、版权/OA 选择和最终归档。
\end{itemize}
\section{1. 撰写与定稿}
\textbf{阶段目标：}以中文母稿为内容基线，固定科学内容、格式和交叉引用后，再进入中文期刊适配或英文翻译。
\subsection{中文母稿}
\begin{itemize}
\item 正文内容、章节顺序、符号定义和结论已经定稿。
\item 所有公式编号一致；Word 稿中的公式全部使用 MathType 编辑。
\item 公式、章节、图、表和参考文献均使用交叉引用。
\item 每张图表均在正文中被引用、描述并解释其支撑的结论。
\item 作者姓名、单位、顺序、通讯作者和基金信息已核对。
\end{itemize}
\subsection{中英文分支}
\begin{itemize}
\item 中文期刊跳过全文翻译，按目标期刊模板整理，并单独完成所需英文题目、摘要、关键词和 Cover Letter。
\item 英文期刊忠实翻译，不增加或删减科学内容，不改变公式、数值、图表和引用。
\item 优先使用计算数学、计算物理、CMAME、JCP、PINNs 和 AI4Science 领域常用术语，避免生僻词与复杂长句。
\item 线下逐句校稿开始后，以英文稿作为提交主版本，中文稿仅用于对照原意。
\end{itemize}
\subsection{目标期刊预演}
\begin{itemize}
\item 提前完整走通投稿系统，以清晰标记的占位文件完成可填写步骤。
\item 预先录入作者、单位、邮箱、ORCID、基金、声明和推荐/回避审稿人信息。
\item 记录文章类型、篇幅、图表、补充材料、数据、代码和文件命名要求。
\item 正式提交前替换全部占位文件。
\end{itemize}
\subsection{线下校稿与 AI 核查}
\begin{itemize}
\item 准备中文稿、英文稿以及中英文 Cover Letter 初稿。
\item 团队逐字逐句审阅英文稿，核对科学含义、术语、逻辑、语法和期刊格式。
\item 团队修改完成后，使用标注式 AI 提示词核查语法与格式；AI 不得直接覆盖原稿。
\end{itemize}
\section{2. 首次投稿}
\textbf{阶段目标：}锁定唯一正式版本，完成系统上传、最终生成稿核对、归档和汇报。
\subsection{Cover Letter}
\begin{itemize}
\item 使用统一开头、结尾和声明框架，只替换论文内容、贡献、期刊匹配度和必要说明。
\item 简洁说明科学问题、方法、主要贡献和期刊匹配度，不复制摘要，不夸大创新。
\end{itemize}
\subsection{版本与上传}
\begin{itemize}
\item 建立唯一待正式提交目录，清除旧稿、临时稿和占位文件。
\item 逐项核对主稿、Cover Letter、图表、补充材料、声明、数据和代码文件。
\item 系统作者、单位、基金、摘要和关键词与主稿一致。
\item 打开期刊生成 PDF，从首页至参考文献逐页核对后再提交。
\end{itemize}
\subsection{归档与汇报}
\begin{itemize}
\item 目录名：正式提交版\_期刊简称\_YYYY-MM-DD。
\item 邮件主题：期刊简称\_正式提交版\_论文短标题。
\item 压缩包发送刘老师邮箱；期刊生成 PDF 优先，否则发送实际上传 PDF，并单独发微信。
\end{itemize}
\section{3. 返修与逐点回复}
\textbf{阶段目标：}先拆分并翻译决定信和审稿意见，再同步修改稿件与 Response to Reviewers。
\subsection{意见整理}
\begin{itemize}
\item 完整保留编辑和审稿意见原文，制作中英文对照版。
\item 编辑意见使用 E.1、E.2；审稿意见使用 R1.1、R1.2、R2.1 等稳定编号。
\item 标明拟采取的修改、待补信息、正文位置和当前状态。
\end{itemize}
\subsection{线下讨论}
\begin{itemize}
\item 准备修改稿和返修信的中英文混排版本，或分别准备中英文文件，以便逐句对照。
\item 工作稿保留 Reviewer comment、中文理解、修改方案、英文回复、正文位置和待补信息链条。
\end{itemize}
\subsection{逐点回复}
\begin{itemize}
\item 每条回复先直接回答，再说明具体修改、科学依据和章节/页码/行号/图表位置。
\item 新增或修改图表时，在返修信中同步展示。
\item 不能完成的建议以研究范围、证据边界或设计限制专业说明，不以时间或方便为主要理由。
\item 不虚构实验、数据、引用、行号、图表或已完成修改；缺失项标记 AUTHOR\_INPUT\_NEEDED。
\end{itemize}
\subsection{返修归档}
\begin{itemize}
\item 目录名：返修提交版\_期刊简称\_ManuscriptID\_R1\_YYYY-MM-DD。
\item 包含决定信、意见双语版、修改前后稿、修订痕迹稿、返修信、上传文件和系统生成 PDF。
\item 邮件和微信同步规则与首次投稿一致，并注明返修轮次和稿件编号。
\end{itemize}
\section{4. 录用与出版前校样}
\textbf{阶段目标：}逐项响应编辑和出版部门要求，保护最终出版信息与科学内容的一致性。
\subsection{Proof 核查}
\begin{itemize}
\item 逐条处理编辑、排版和版权系统问题并提供所需素材。
\item 核对标题、作者、单位、通讯信息、基金、公式、图表、参考文献和补充材料。
\item 校样阶段原则上只纠正排版、文字和事实错误，不无说明地改变核心结论。
\item 提交前保存最终 Proof、问题回复、版权协议和系统确认页。
\end{itemize}
\subsection{OA 与订阅}
\begin{itemize}
\item 默认优先订阅模式，不主动选择 OA。
\item 只有期刊政策、基金要求或刘老师明确确认时选择 OA，并核对费用和许可。
\end{itemize}
\section{命名、归档与通知}
\begin{itemize}
\item 首次投稿目录：\texttt{正式提交版\_期刊简称\_YYYY-MM-DD}
\item 返修目录：\texttt{返修提交版\_期刊简称\_ManuscriptID\_R1\_YYYY-MM-DD}
\item 首次投稿邮件：\texttt{期刊简称\_正式提交版\_论文短标题}
\item 返修邮件：\texttt{期刊简称\_返修提交版\_R1\_论文短标题}
\end{itemize}
\section{中译英提示词}
\begin{verbatim}
你是一名熟悉计算数学、计算物理、CMAME、Journal of Computational Physics、PINNs、深度学习求解微分方程与 AI4Science 写作规范的学术翻译编辑。

任务：将下面的中文论文内容翻译为清晰、准确、自然的学术英文。

必须遵守：
1. 不改变原意，不增加或删减科学内容，不强化或弱化结论。
2. 保留所有公式、符号、变量、数值、单位、图表编号、章节编号和参考文献标记。
3. 优先使用上述领域和目标期刊常见术语，避免生僻词、口语、过度华丽表达与过长句。
4. 同一术语、方法名和缩写全文一致；首次出现的缩写按学术规范展开。
5. 若中文原文存在歧义、逻辑跳跃或术语不确定，不要擅自推断；在译文后列出“需要作者确认的问题”。
6. 只输出：A. 英文译文；B. 术语对照表；C. 需要作者确认的问题（若无则写 None）。

目标期刊：[JOURNAL]
文章部分：[ABSTRACT / INTRODUCTION / METHODS / RESULTS / DISCUSSION]
术语表或既定译法：[GLOSSARY]

待翻译内容：
[PASTE CHINESE TEXT]
\end{verbatim}
\section{英文语法与格式核查提示词}
\begin{verbatim}
你是一名计算数学、计算物理与 AI4Science 论文的英文校稿员。请对下面文本执行“标注式核查”，不要直接重写全文。

核查范围：语法、拼写、标点、冠词、时态、主谓一致、学术用语、术语一致性、公式/图表/参考文献交叉引用、期刊格式和可读性。

红线：
1. 不改变科学含义、数据、公式、数值、单位、结论或引用。
2. 不虚构缺失内容，不把相关性改成因果性，不夸大创新性。
3. 对每个问题给出位置、原文、问题类型、建议修改和理由。
4. 将可能改变原意的建议标记为“需作者确认”，不得直接采用。

输出格式：
| 编号 | 位置/原文 | 问题类型 | 建议修改 | 理由 | 风险等级 |

最后给出：术语一致性清单、交叉引用问题清单、期刊格式问题清单，以及“是否存在可能改变科学含义的修改”。

目标期刊及格式要求：[JOURNAL AND GUIDELINES]
待核查文本：
[PASTE ENGLISH TEXT]
\end{verbatim}
\section{Cover Letter 模板}
\begin{verbatim}
Dear [Editor-in-Chief / Handling Editor],

Please find enclosed our manuscript entitled “[MANUSCRIPT TITLE]”, which we wish to submit for consideration as a [ARTICLE TYPE] in [JOURNAL NAME].

This work addresses [SCIENTIFIC PROBLEM]. We develop [METHOD / FRAMEWORK] to [PRIMARY PURPOSE]. The main contributions are as follows: (1) [CONTRIBUTION 1]; (2) [CONTRIBUTION 2]; and (3) [CONTRIBUTION 3]. The results demonstrate [MAIN RESULT WITHOUT OVERSTATEMENT].

We believe that the manuscript is well suited to [JOURNAL NAME] because [SCOPE AND READERSHIP FIT]. It may be of particular interest to readers working on [RELEVANT TOPICS].

This manuscript is original, has not been published previously, and is not under consideration for publication elsewhere. All authors have approved the manuscript and agree with its submission to [JOURNAL NAME]. [CONFLICT OF INTEREST / DATA / CODE STATEMENT IF REQUIRED.]

Thank you for considering our manuscript. We look forward to your response.

Sincerely,
[CORRESPONDING AUTHOR]
[AFFILIATION]
[POSTAL ADDRESS]
[EMAIL]
[ORCID, IF REQUIRED]
\end{verbatim}
\section{Response to Reviewers 模板}
\begin{verbatim}
Dear Editor and Reviewers,

We thank the editor and reviewers for their careful evaluation of our manuscript entitled “[TITLE]” (Manuscript ID: [ID]). We have revised the manuscript to address the concerns raised and provide a point-by-point response below. Reviewer comments are reproduced in full, followed by our responses and the corresponding manuscript changes.

Response strategy summary
- Decision type: [MAJOR / MINOR REVISION]
- Package readiness: [draft_with_placeholders / needs_author_input / ready_to_submit]
- Major risks or missing information: [ITEMS / NONE]

Editor comment E.1
[FULL EDITOR COMMENT]

Response
[DIRECT ANSWER.] To address this point, we have [SPECIFIC ACTION]. The corresponding change appears in [SECTION, PAGE, LINES, FIGURE OR TABLE].

Reviewer 1 comment R1.1
[FULL REVIEWER COMMENT]

Response
We thank the reviewer for raising this point. [DIRECT ANSWER.] We have [SPECIFIC REVISION OR ANALYSIS]. This change appears in [SECTION, PAGE, LINES, FIGURE OR TABLE]. [REMAINING LIMITATION, IF ANY.]

中文核对
- 对该意见的中文理解：[CHINESE INTERPRETATION]
- 实际完成的修改：[AUTHOR-SUPPLIED ACTION]
- 待补事实或位置：[AUTHOR_INPUT_NEEDED / NONE]

Manuscript change checklist
- E.1: [ACTION AND LOCATION]
- R1.1: [ACTION AND LOCATION]

Sincerely,
The Authors
\end{verbatim}
\section{邮件模板}
\begin{verbatim}
主题：[期刊简称]_[正式提交版 / 返修提交版_R1]_[论文短标题]

刘老师您好：

论文《[论文题目]》已于 [YYYY-MM-DD] 完成向 [期刊全名] 的[首次投稿 / 第 R1 轮返修提交]。

关键信息：
- 期刊：[期刊全名]
- 稿件编号：[Manuscript ID / 尚未生成]
- 提交日期：[YYYY-MM-DD]
- 提交类型：[首次投稿 / Major Revision / Minor Revision]
- 最终 PDF：[期刊生成版本 / 实际上传版本]

附件为本次完整提交文件压缩包，包含中英文稿、Cover Letter / Response to Reviewers、全部上传文件及最终生成版本。

请您查收。
\end{verbatim}
\section{微信模板}
\begin{verbatim}
刘老师您好，《[论文短标题]》已于 [YYYY-MM-DD] 完成向 [期刊简称] 的[正式投稿 / 第 R1 轮返修提交]，稿件编号为 [ID / 尚未生成]。现将最终 PDF 发您查阅；完整提交材料已发送至您的邮箱。
\end{verbatim}
\end{document}